\documentclass[a4paper,11pt]{article}

% Packages
\usepackage[utf8]{inputenc}
\usepackage[T1]{fontenc}
\usepackage[french]{babel}
\usepackage{amsmath, amssymb}
\usepackage{geometry}
\usepackage{xcolor}
\usepackage{enumitem}
\usepackage{fancyhdr}
\usepackage{array}
\usepackage{booktabs}
\usepackage{graphicx}
\usepackage{titlesec}
\usepackage[
    colorlinks=true,
    linkcolor=black,
    urlcolor=black
]{hyperref}
\usepackage{bookmark}  % ← corrige le conflit titlesec/hyperref

% Mise en page
\geometry{margin=2.5cm}
\setlength{\headheight}{14pt}
\pagestyle{fancy}
\fancyhf{}
\fancyhead[L]{\textit{Assistant Intelligent pour les Cours de L1}}
\fancyhead[R]{\textit{CMI Informatique --- 2025/2026}}
\fancyfoot[C]{\thepage}

% Style des sections
\titleformat{\section}
    {\large\bfseries}
    {\thesection.}{0.6em}{}
    [\vspace{2pt}\titlerule\vspace{4pt}]

\titleformat{\subsection}
    {\normalsize\bfseries}
    {\thesubsection.}{0.6em}{}

% ============================================================
\begin{document}

% ============================================================
% PAGE DE COUVERTURE
% ============================================================
\thispagestyle{empty}

\vspace*{3cm}

\begin{center}
    {\small\scshape Projet de Recherche --- CMI Informatique}\\[1cm]
    {\rule{\linewidth}{0.4pt}}\\[0.6cm]
    {\LARGE\bfseries Création d'un Assistant Intelligent\\[0.3cm]pour les Cours de Première Année}\\[0.6cm]
    {\rule{\linewidth}{0.4pt}}\\[1.2cm]
    {\normalsize\itshape Public cible : Tous}\\[0.4cm]
    {\normalsize Encadrant : Yannick Estève}\\[2.5cm]
    {\normalsize
        A. Louvet \quad L. Ouelhadj \quad T. De Faveri-Serre \quad C. Cloupet \quad H. Baudey
    }\\[0.8cm]
    {\small Années universitaires 2025--2027}
\end{center}

\vfill

\begin{center}
\noindent\rule{0.5\linewidth}{0.4pt}\\[0.5cm]
{\small
Ce projet vise d'une part à produire des contributions de recherche sur les
grands modèles de langage (LLM), la génération augmentée par récupération
(RAG), la recherche vectorielle et les frameworks d'orchestration dans le
contexte des assistants pédagogiques, en les rendant accessibles quelle que
soit la familiarité préalable avec ces concepts ; et d'autre part à
construire un assistant conversationnel ancré dans le RAG, capable de
répondre aux questions des étudiants à partir de leurs propres documents de
cours.
}\\[0.4cm]
\noindent\rule{0.5\linewidth}{0.4pt}
\end{center}

\vspace{2cm}

% ============================================================
% TABLE DES MATIÈRES
% ============================================================

\newpage
\tableofcontents

% ============================================================
% PARTIE 1 : ÉTAT DE L'ART (DÉFINITIONS ET CONTEXTE)
% ============================================================

\newpage
\section{État de l'Art : Définitions et Contexte}

\subsection{Introduction à l'Intelligence Artificielle}

\subsubsection{Bref historique}
L'idée de machines capables de penser n'est pas nouvelle. Dès 1950, Alan
Turing proposa son célèbre \emph{Test de Turing} comme critère d'intelligence
machine \cite{turing1950computing} : si une machine peut soutenir une
conversation indiscernable de celle d'un être humain, on peut, dans un certain
sens, dire qu'elle pense. Le terme « Intelligence Artificielle » lui-même fut
inventé en 1956 par John McCarthy lors de la conférence de Dartmouth,
considérée comme l'acte fondateur du domaine.

Les décennies suivantes alternèrent entre périodes d'optimisme et
« hivers de l'IA » --- des intervalles de réduction des financements et de
l'intérêt, causés par l'écart entre les promesses ambitieuses et les résultats
modestes. Le paysage changea radicalement dans les années 1980 et 1990 avec
l'essor de l'\emph{apprentissage automatique} : plutôt que de coder des règles
en dur dans un programme, les chercheurs commencèrent à concevoir des systèmes
capables d'\emph{apprendre} des patterns directement à partir des données.
Ce changement de paradigme posa les fondements de tout ce qui suivit.

Le véritable tournant arriva en 2012, quand un réseau de neurones profond
appelé AlexNet \cite{krizhevsky2012alexnet} écrasa la concurrence lors du défi
de reconnaissance visuelle ImageNet, démontrant que l'\emph{apprentissage
profond} --- des réseaux de neurones à nombreuses couches entraînés sur de
grands jeux de données --- pouvait atteindre des performances bien supérieures
aux méthodes classiques. À partir de ce moment, l'apprentissage profond
devint l'approche dominante dans presque tous les sous-domaines de l'IA.

\subsubsection{Qu'est-ce que l'Intelligence Artificielle ?}
Le terme \emph{Intelligence Artificielle} est l'une des expressions les plus
utilisées --- et les plus galvaudées --- du paysage technologique actuel. Avant
d'aborder les systèmes spécifiques au cœur de ce projet, il convient de
prendre du recul pour comprendre ce qu'est réellement l'IA, d'où elle vient
et pourquoi elle est devenue si centrale en informatique moderne.

Dans son essence, l'Intelligence Artificielle désigne tout système de calcul
conçu pour effectuer des tâches qui nécessiteraient normalement l'intelligence
humaine --- comme reconnaître des objets dans une image, comprendre une phrase
prononcée, ou décider du prochain coup dans une partie d'échecs. Une définition
plus précise, proposée par Russell et Norvig dans leur manuel de référence
\cite{russell2010ai}, décrit un agent IA comme un système qui \emph{perçoit
son environnement à travers des capteurs et agit sur lui par des actionneurs
afin de maximiser une mesure de performance}.

Une distinction utile est celle entre l'\emph{IA étroite} (ou IA faible) et
l'\emph{IA générale} (ou IA forte). L'IA étroite --- qui recouvre quasiment
tous les systèmes existants aujourd'hui --- est conçue pour exceller dans un
type de tâche précis : un filtre anti-spam classe les e-mails, un moteur de
recommandation suggère des films, un modèle de langage génère du texte.
L'IA générale, en revanche, serait un système capable de raisonner sur des
domaines arbitraires au niveau humain ; elle reste largement hypothétique et
constitue un sujet de recherche à long terme. Tous les systèmes d'IA évoqués
dans ce projet relèvent de l'IA étroite.

% ------------------------------------------------------------
\newpage
\subsection{Les Grands Modèles de Langage (LLM)}

Les grands modèles de langage --- LLM, pour \emph{Large Language Models}, car
le nom complet est un peu long --- sont bien plus présents dans notre quotidien
qu'on ne le pense, ou même qu'on ne l'anticipait il y a encore quelques années.
Mais que sont-ils vraiment ? Que recouvrent-ils, et à quoi servent-ils ?

Un grand modèle de langage est, pour faire simple, un modèle mathématique qui
génère du texte en prédisant le mot le plus probable dans une suite, un peu
comme les « suggestions de mots » qui apparaissent en haut du clavier quand
vous tapez un message sur votre téléphone. Mais ne vous y trompez pas : là où
votre smartphone peine à prédire autre chose que votre commande de café
habituelle, un LLM opère à une échelle franchement vertigineuse. Imaginez que
l'on prenne chaque livre, chaque article, chaque ligne de code disponible sur
internet, et qu'on les injecte dans un cerveau numérique qui cartographie les
relations complexes entre chaque mot et chaque concept. En calculant la
probabilité de la suite d'un texte à partir de milliards de paramètres, ces
modèles dépassent largement la simple « saisie automatique » pour entrer dans
le domaine du raisonnement et de la créativité : ils ne devinent pas des
lettres, ils saisissent les nuances de l'intention humaine.

\subsubsection{L'architecture Transformer}

Au cœur de ce processus se trouve une structure appelée \emph{Transformer}
\cite{vaswani2017attention}. Imaginez-la comme un système qui permet au modèle
d'examiner un paragraphe entier d'un seul coup, plutôt qu'un mot à la fois.
Ce mécanisme d'« \emph{attention} » garantit que le modèle se souvient du
début d'un texte pendant qu'il en rédige la fin, maintenant un fil conducteur
cohérent qui paraît remarquablement humain.

\subsubsection{Un outil polyvalent}

Parce qu'ils sont des maîtres de la reconnaissance de motifs (\emph{pattern
recognition}), leur utilité s'étend bien au-delà de la réponse à des questions
triviales ou de la rédaction d'e-mails :

\begin{itemize}
  \item \textbf{Partenaires créatifs :} Ils peuvent générer des slogans
    marketing ou rédiger de la poésie dans le style d'un romancier du
    XIX\textsuperscript{e} siècle.
  \item \textbf{Assistants techniques :} Ils traduisent des langages de
    programmation complexes aussi facilement qu'ils traduisent le français en
    anglais.
  \item \textbf{Synthétiseurs de données :} Ils peuvent résumer un document
    juridique de 50 pages en trois points en un clin d'œil.
\end{itemize}

En essence, un LLM agit comme un pont entre le monde rigide et binaire des
ordinateurs et le monde désordonné et fluide du langage humain. Ils ne
« savent » pas les choses comme nous les savons, mais ils sont des imitateurs
de premier ordre de la logique, ce qui en fait l'outil le plus polyvalent de
l'ère numérique. À mesure que ces modèles continuent d'être perfectionnés, ils
passent de simples curiosités à de véritables collaborateurs essentiels,
transformant non seulement la façon dont nous écrivons, mais aussi la façon
dont nous résolvons des problèmes et interagissons avec la connaissance
collective du monde.

\newpage
\subsubsection{Limites : Hallucinations et Frontières de la Connaissance}

Malgré leurs performances impressionnantes, les LLM souffrent de plusieurs
limitations bien documentées. La plus critique dans un contexte pédagogique
est la tendance à produire des \emph{hallucinations} --- générer du texte
fluide et plausible mais factuellement incorrect ou sans fondement
\cite{zhang2023hallucination, huang2025hallucination}.

Les hallucinations surviennent pour plusieurs raisons. Premièrement, les LLM
reposent sur des corrélations statistiques dans leurs données d'entraînement
plutôt que sur un raisonnement factuel ancré. Deuxièmement, leur connaissance
est statique : ils ont une date limite d'entraînement et ne peuvent pas
accéder à des informations nouvelles ou spécifiques à un domaine sauf si elles
leur sont explicitement fournies. Troisièmement, lorsqu'une requête sort de la
zone de connaissance assurée du modèle, celui-ci peut fabriquer une réponse
plutôt qu'admettre son incertitude. Selon une enquête de 2024 publiée dans
l'\emph{ACM Transactions on Information Systems}, les hallucinations sont
catégorisées en deux types principaux : l'\emph{hallucination de facticité}
(énoncer des faits incorrects) et l'\emph{hallucination de fidélité} (produire
du contenu incohérent avec une source donnée).

Ces limitations sont particulièrement problématiques dans un contexte
pédagogique où l'exactitude est essentielle. Un assistant qui fournit en toute
confiance des informations erronées sur un concept de cours pourrait activement
induire les étudiants en erreur. C'est ce qui motive le recours à l'ancrage
documentaire externe via le RAG.

\subsubsection{Biais et Dégradation par Boucle de Rétroaction via le RLHF}
Une limitation plus subtile mais significative concerne le mécanisme de
retour d'information utilisé pour améliorer les LLM. Des techniques comme
l'\emph{Apprentissage par Renforcement à partir des Retours Humains} (RLHF)
et l'IA Constitutionnelle sont employées pour aligner le comportement du
modèle avec les préférences humaines et réduire les sorties nuisibles. Ce
système repose sur un principe dans lequel les utilisateurs notent positivement
ou négativement les réponses du modèle, qui est ensuite mis à jour en
conséquence \cite{ouyang2022rlhf}. Bien que ce processus vise à aligner le
comportement du modèle avec les attentes des utilisateurs, il peut produire
des effets pervers lorsque ceux-ci soumettent systématiquement des requêtes
mal formulées.

Lorsqu'un utilisateur pose une question ambiguë, incomplète ou incohérente,
le modèle n'a pas de base fiable sur laquelle produire une réponse pertinente.
Il génère alors une réponse statistiquement plausible au regard de l'entrée,
mais qui ne correspond inévitablement pas à l'intention non exprimée de
l'utilisateur. Celui-ci, percevant la réponse comme hors sujet ou absurde,
attribue une note négative --- alors que, du point de vue du modèle, la
réponse était une complétion raisonnable de la requête donnée. Ce désalignement
entre le signal (une note négative) et sa cause réelle (une requête mal
formulée plutôt qu'une erreur du modèle) introduit une \emph{supervision
corrompue} dans la boucle d'entraînement.

Répétée à grande échelle, cette dynamique pousse le modèle à associer certains
patterns de réponse légitimes à une récompense négative et à les désapprendre,
entraînant une dégradation mesurable de la qualité générale des réponses au
fil du temps. Une étude de Chen et al.\ \cite{chen2023chatgpt} a documenté
empiriquement ce phénomène pour ChatGPT, observant des changements de
comportement significatifs entre les versions du modèle, avec certaines tâches
où les versions antérieures surpassaient les versions ultérieures. Les auteurs
attribuent en partie cette régression aux effets cumulatifs du bruit dans les
retours d'information des utilisateurs dans le pipeline RLHF.

Ce problème est particulièrement pertinent dans un contexte pédagogique, où
les étudiants --- surtout en première année --- peuvent manquer de compétences
en formulation de requêtes pour poser des questions précises. Un assistant qui
se dégrade sous des retours bruités pourrait ainsi devenir progressivement
moins utile pour son public principal, soulignant l'importance de concevoir
des mécanismes robustes de collecte de retours et, dans la mesure du possible,
de réduire la dépendance aux évaluations non filtrées des utilisateurs pour
les mises à jour du modèle.

% ------------------------------------------------------------
\newpage
\subsection{Récapitulatif}
Un point de départ naturel pour cet aperçu est une question qui semble simple
mais s'avère en réalité assez nuancée : qu'est-ce exactement qu'un LLM et en
quoi diffère-t-il de la notion plus large d'« IA » ? Le terme \emph{Intelligence
Artificielle} (IA) est un concept chapeau englobant toute technique ou système
permettant aux machines de reproduire l'intelligence humaine, incluant
l'apprentissage automatique, la vision par ordinateur, la robotique et le
traitement du langage naturel. Un \emph{Grand Modèle de Langage} (LLM), en
revanche, désigne une catégorie spécifique de modèle IA conçu pour comprendre
et générer le langage humain. Les LLM sont un sous-ensemble de l'IA, au même
titre que le \emph{Machine Learning} (ML) ou le \emph{Deep Learning} (DL),
etc\dots ; si tout LLM est un système IA, la réciproque n'est pas vraie.

Concrètement, les LLM sont des réseaux de neurones entraînés sur d'immenses
corpus de texte par apprentissage auto-supervisé, fondés sur l'architecture
Transformer introduite par Vaswani et al.\ \cite{vaswani2017attention}, qui
permet aux modèles de capturer les dépendances à longue portée dans le texte
grâce aux mécanismes d'auto-attention. L'introduction ultérieure de la série
GPT d'OpenAI --- GPT-2 en 2019, GPT-3 en 2020, puis GPT-4 --- a démontré
que l'augmentation de la taille du modèle et des données d'entraînement
conduit à des \emph{capacités émergentes} en compréhension et génération du
langage \cite{wei2022emergent} : des capacités qui n'apparaissent
qualitativement qu'au-delà d'un certain seuil de taille. Les LLM contemporains
tels que GPT-4, Claude d'Anthropic, LLaMA de Meta, Gemini de Google et
Mistral font preuve d'une remarquable maîtrise dans un large éventail de
tâches --- de la réponse aux questions et la synthèse à la génération de
code et au raisonnement multi-étapes --- et se caractérisent par des nombres
de paramètres de l'ordre de dizaines ou centaines de milliards.

% ------------------------------------------------------------
\subsection{Génération Augmentée par Récupération (RAG)}

\subsubsection{Principe et Motivation}
Avant de plonger dans l'architecture technique, il vaut la peine de prendre
du recul pour expliquer ce qu'est le RAG et pourquoi il a été développé,
afin de rendre cette recherche accessible et intéressante pour tous.

Imaginez demander à un ami très cultivé une question sur un sujet précis ---
disons, le contenu d'un cours auquel il n'a jamais assisté. Peu importe son
niveau d'intelligence, il ne peut tout simplement pas savoir ce qui a été
enseigné dans ce cours particulier. Un LLM fait face à exactement la même
limitation : il ne connaît que ce qu'il a vu lors de l'entraînement, et n'a
aucun moyen de consulter des documents qu'il n'a jamais rencontrés. Imaginez
maintenant qu'avant de répondre, votre ami soit autorisé à feuilleter
rapidement les notes de cours pertinentes et à utiliser ce qu'il y trouve
pour informer sa réponse. C'est en essence ce que fait le RAG : il donne au
modèle accès à un ensemble de documents spécifiques au moment de répondre,
de sorte que sa réponse est ancrée dans ces sources plutôt que dans sa mémoire
générale (et potentiellement obsolète ou incomplète).

Plus formellement, la Génération Augmentée par Récupération (RAG) est un
paradigme qui améliore les LLM en les connectant à une base de connaissances
externe au moment de l'inférence \cite{lewis2020rag}. Plutôt que de s'appuyer
uniquement sur la mémoire paramétrique du modèle, un système RAG récupère
d'abord les passages de documents pertinents d'un corpus curé en fonction de
la requête de l'utilisateur, puis transmet ces passages comme contexte au LLM,
qui génère une réponse ancrée.

Cette approche répond directement aux limitations décrites dans la section
précédente : elle fournit au LLM des informations à jour et spécifiques au
domaine, et ancre la génération dans des sources explicites et vérifiables
--- réduisant significativement le risque d'hallucination. Une enquête
spécifiquement focalisée sur le RAG pour les applications pédagogiques
\cite{rag_education_survey} confirme que cette approche améliore
substantiellement la pertinence et l'exactitude factuelle des réponses
générées par l'IA dans les contextes académiques.

\newpage
\subsubsection{Architecture du RAG}

Un pipeline RAG standard comprend trois étapes :

\begin{enumerate}
  \item \textbf{Indexation (phase développeur) :} Les documents sources (PDF,
    diaporamas de cours, notes de lecture, etc.) sont découpés en morceaux
    (\emph{chunks}), encodés sous forme de vecteurs d'embedding à l'aide d'un
    modèle d'embedding, et stockés dans une base de données vectorielle.
    \emph{(cf. Section 1.4 avec les détails sur la recherche vectorielle et
    l'indexation)}.

  \item \textbf{Récupération (phase utilisateur) :} Lorsqu'un utilisateur
    soumet une requête, celle-ci est encodée dans le même espace d'embedding.
    Le système effectue ensuite une recherche par similarité pour récupérer
    les $k$ morceaux les plus pertinents de la base de données.

  \item \textbf{Génération :} Les morceaux récupérés sont ajoutés en préfixe
    à la requête de l'utilisateur comme contexte dans le prompt du LLM. Le
    LLM génère ensuite une réponse ancrée dans ces informations récupérées.
\end{enumerate}

\subsubsection{Variantes et Évolution du RAG}

Le domaine a considérablement mûri depuis l'article original sur le RAG.
Des enquêtes de 2024--2025 identifient trois paradigmes principaux
\cite{gao2024rag} :

\begin{itemize}
  \item \textbf{RAG Naïf :} Le pipeline de base récupération-puis-génération
    décrit ci-dessus.
  \item \textbf{RAG Avancé :} Introduit la réécriture de requêtes, le
    re-classement des résultats récupérés et la récupération itérative pour
    améliorer la précision.
  \item \textbf{RAG Modulaire :} Une architecture plus flexible où chaque
    composant (récupérateur, lecteur, mémoire, etc.) peut être échangé ou
    optimisé indépendamment.
\end{itemize}

Des travaux plus récents explorent également le \textbf{RAG Agentique}, où le
modèle peut décider \emph{quand} récupérer et \emph{quelles} sources consulter
en fonction de la complexité de la requête, en utilisant une boucle de
raisonnement multi-étapes \cite{zhao2024rag}.

Pour notre projet, une architecture RAG Naïf ou Avancé est suffisante étant
donné le périmètre bien défini et borné du corpus documentaire (supports de
cours de première année).

% ------------------------------------------------------------
\subsection{Recherche Vectorielle et Indexation}

\subsubsection{Qu'est-ce qu'un Embedding ?}

Un \emph{embedding}, dans le contexte de l'intelligence artificielle et des
grands modèles de langage (LLM), désigne la transformation d'un élément
d'information --- un mot, une phrase, un document, ou parfois même une image
--- en une séquence de nombres réels appelée \emph{vecteur}. Cette
représentation numérique permet à la machine de manipuler le langage et les
concepts non plus comme du texte symbolique, mais comme des objets
mathématiques sur lesquels elle peut effectuer des calculs.

Plutôt que de « comprendre » les mots comme le ferait un humain, le modèle
les place dans un \emph{espace vectoriel} de haute dimension (typiquement
entre 384 et 4096 dimensions selon le modèle), où chaque élément occupe une
position précise. Dans cet espace, la \emph{distance} entre deux vecteurs
reflète leur \emph{similarité sémantique} : deux phrases de sens proches
auront des vecteurs voisins, même si elles ne partagent aucun mot en commun.
C'est le principe qui permet à un système, par exemple, de récupérer le
passage \textit{« la photosynthèse permet aux plantes de produire de
l'énergie »} à partir de la requête \textit{« comment les plantes
génèrent-elles leur énergie ? »}.

\newpage
\subsubsection{Comment les Embeddings sont-ils Appris ?}

Ces vecteurs ne sont pas choisis arbitrairement : ils sont \emph{appris}
lors de l'entraînement du modèle. Le système lit des quantités immenses de
texte et ajuste progressivement ses paramètres de sorte que les vecteurs
résultants capturent les régularités du langage. Par exemple, le modèle
apprend que « chat » et « chien » apparaissent dans des contextes similaires
et finit par leur attribuer des vecteurs voisins, tandis que « voiture »
sera plus éloigné.

Des travaux pionniers comme \emph{Word2Vec} \cite{mikolov2013word2vec} et
\emph{GloVe} \cite{pennington2014glove} ont montré que ce simple principe
suffit à faire émerger des relations sémantiques remarquables. L'exemple
canonique d'arithmétique vectorielle,
$\overrightarrow{\text{roi}} - \overrightarrow{\text{homme}}
+ \overrightarrow{\text{femme}} \approx \overrightarrow{\text{reine}}$,
illustre que l'espace appris encode des concepts (genre, royauté) plutôt que
de simples co-occurrences.

Dans les modèles modernes basés sur les Transformers, les embeddings sont
également \emph{contextuels} : le vecteur associé à un mot dépend de la
phrase dans laquelle il apparaît. Le mot « banque » dans \textit{« je me
suis assis au bord de la rivière »} et dans \textit{« j'ai déposé de
l'argent à la banque »} reçoit ainsi deux représentations distinctes, ce qui
résout naturellement les ambiguïtés lexicales. Pour la recherche sémantique,
on s'appuie généralement sur des modèles spécialisés comme ceux de la
famille \texttt{sentence-transformers} \cite{reimers2019sbert}, qui
produisent un \emph{vecteur unique} par phrase ou paragraphe, optimisé pour
la comparaison par similarité plutôt que pour la génération de texte.

\subsubsection{Mesurer la Similarité entre Vecteurs}

Une fois le contenu représenté sous forme de vecteurs, comparer deux textes
revient à comparer deux points dans l'espace. Trois métriques sont
couramment utilisées :

\begin{itemize}
  \item \textbf{La similarité cosinus}, qui mesure l'angle entre deux
    vecteurs et est insensible à leur magnitude : de loin la métrique la
    plus utilisée en recherche sémantique.
  \item \textbf{Le produit scalaire}, plus rapide à calculer mais sensible
    aux normes des vecteurs ; préféré lorsque les vecteurs ont été normalisés
    au préalable.
  \item \textbf{La distance euclidienne}, intuitive mais souvent moins
    discriminante en très haute dimension en raison du « fléau de la
    dimensionnalité ».
\end{itemize}

\subsubsection{Indexation et Recherche Approchée du Plus Proche Voisin}

Lorsqu'un corpus contient des milliers, voire des millions de passages
indexés, comparer la requête à chaque vecteur stocké --- la recherche dite
\emph{exacte} ou « force brute » --- devient prohibitivement coûteux. La
solution standard consiste à s'appuyer sur des structures d'\emph{indexation}
permettant la \emph{recherche approchée du plus proche voisin} (ANN) : on
échange une petite perte de précision contre plusieurs ordres de grandeur en
rapidité.

Plusieurs algorithmes d'indexation sont devenus des standards :

\begin{itemize}
  \item \textbf{IVF (Index de Fichier Inversé)} : partitionne l'espace en
    cellules via une étape préalable de regroupement (typiquement $k$-means),
    puis restreint la recherche aux cellules les plus proches de la requête.
  \item \textbf{HNSW (Monde de Petits Mondes Hiérarchique Navigable)}
    \cite{malkov2018hnsw} : construit un graphe multi-couches dans lequel la
    navigation locale converge rapidement vers le voisinage du vecteur
    requêté. C'est actuellement l'algorithme de référence dans la plupart des
    bases de données vectorielles, grâce à son excellent compromis entre
    rappel et latence.
  \item \textbf{Quantification par Produit (PQ)} : compresse les vecteurs
    en sous-quantifiant chaque sous-espace, réduisant drastiquement
    l'utilisation mémoire au prix d'une légère perte de précision.
\end{itemize}

Ces algorithmes sont implémentés dans des bibliothèques et bases de données
dédiées --- les \emph{bases de données vectorielles} --- les plus répandues
dans l'écosystème open-source étant \texttt{FAISS} (Facebook AI Similarity
Search), \texttt{Chroma}, \texttt{Qdrant} et \texttt{Weaviate}. Le choix
dépend principalement du volume de données, des contraintes de latence et de
la nécessité d'un filtrage par métadonnées (par exemple, restreindre la
recherche à un cours ou un chapitre spécifique).

\subsubsection{Synthèse}

En pratique, les embeddings constituent le bloc de construction fondamental
de tout système RAG : ils sont utilisés pour interpréter les requêtes,
comparer les passages de cours, effectuer la recherche sémantique et
alimenter plus largement les mécanismes internes des LLM. On peut les voir
comme une « traduction » du langage humain en une géométrie de concepts, où
les idées sont organisées dans un espace multidimensionnel que la machine
peut explorer et exploiter efficacement. C'est précisément cette géométrie,
combinée à un index adapté, qui permettra à notre assistant pédagogique de
récupérer en quelques millisecondes les passages de cours pertinents pour
répondre à la question d'un étudiant.

% ============================================================
% PARTIE 2 : PRÉSENTATION DU PROJET
% ============================================================
\newpage
\section{Présentation du Projet}

\subsection{Objectifs Pédagogiques}
L'objectif est de découvrir les principes fondamentaux de l'intelligence
artificielle moderne à travers les grands modèles de langage (LLM) et le
système RAG (Génération Augmentée par Récupération). Les étudiants apprendront
à interfacer un LLM avec une base documentaire personnalisée pour créer un
assistant conversationnel capable de répondre à des questions portant sur
leurs cours.

\subsection{Les Assistants IA en Contexte Éducatif}

\subsubsection{Panorama des Travaux Existants}

L'application d'assistants IA à l'éducation suscite un intérêt de recherche
croissant. Une enquête publiée dans \emph{Applied Sciences}
\cite{mdpi2025rag} a recensé 47 études sur les chatbots éducatifs basés sur
le RAG, la majorité publiées en 2024, témoignant de la croissance rapide de
ce domaine.

Plusieurs déploiements réels sont directement pertinents pour notre projet :

\begin{itemize}
  \item \textbf{UniMiB Assistant} (Université de Bicocca, Milan, 2024)
    \cite{unimib2024} : Un assistant étudiant basé sur le RAG, capable de
    répondre aux questions sur les procédures universitaires, les emplois du
    temps et les informations administratives. Il indexe les documents
    officiels de l'université et utilise un LLM pour générer des réponses
    ancrées dans ces sources.

  \item \textbf{OwlMentor} (Frontiers in Psychology, 2024)
    \cite{owlmentor2024} : Un environnement d'apprentissage alimenté par l'IA
    conçu pour aider les étudiants universitaires à comprendre des textes
    scientifiques. Intégré dans un cours, il offrait un chat basé sur les
    documents, la génération automatique de questions et la création de quiz.

  \item \textbf{Assistant IA du CS50} (Université Harvard) : Intégré dans un
    cours d'introduction à l'informatique, un assistant basé sur le RAG s'est
    avéré capable de fournir des réponses détaillées et contextuellement
    appropriées, facilitant le tutorat personnalisé à grande échelle.

  \item \textbf{Assistant RAG pour l'enseignement médical} (Nature
    \emph{npj Digital Medicine}, 2025) \cite{npj2025} : Un assistant
    pédagogique déployé dans un cours de médecine, où les modèles d'utilisation
    et les retours des étudiants ont été analysés sur deux cohortes
    consécutives, démontrant la faisabilité et l'efficacité des assistants
    basés sur le RAG dans des contextes éducatifs à enjeux élevés.
\end{itemize}

\subsubsection{Efficacité et Défis}

Le constat général à travers ces études est que les assistants basés sur le
RAG sont efficaces pour fournir des réponses précises et ancrées dans le
contexte à des questions spécifiques à un domaine, et que les étudiants les
perçoivent positivement comme un complément à l'enseignement humain.
Cependant, plusieurs défis sont régulièrement rapportés :

\begin{itemize}
  \item \textbf{Qualité et couverture des documents :} La qualité des réponses
    de l'assistant est directement bornée par la qualité et la couverture du
    corpus documentaire indexé. Des documents de cours mal formatés ou
    incomplets dégradent les performances.
  \item \textbf{Requêtes hors domaine :} Les étudiants posent souvent des
    questions en dehors du domaine des documents indexés. Les systèmes robustes
    doivent détecter ces requêtes et refuser gracieusement de spéculer.
  \item \textbf{Difficulté d'évaluation :} Évaluer la correction des réponses
    générées dans des contextes éducatifs ouverts est non trivial et nécessite
    souvent une évaluation par des experts humains.
  \item \textbf{Calibration de la confiance des étudiants :} Les étudiants
    peuvent sur-dépendre des réponses de l'IA ou, à l'inverse, s'en méfier
    même lorsqu'elles sont correctes. Communiquer l'incertitude et citer les
    sources est important pour une calibration appropriée de la confiance.
\end{itemize}

% ------------------------------------------------------------
\newpage
\subsection{Synthèse et Positionnement de Notre Projet}

Cet état de l'art révèle que la combinaison des LLM et du RAG représente la
meilleure pratique actuelle pour construire des assistants conversationnels
spécifiques à un domaine. La pile technologique que nous prévoyons d'utiliser
--- Python, un modèle d'embedding (\texttt{sentence-transformers} ou OpenAI),
une base de données vectorielle (FAISS ou Chroma) et un framework
d'orchestration (LangChain ou LlamaIndex) --- est bien validée tant dans la
littérature académique que dans les déploiements éducatifs réels.

L'originalité principale de notre projet réside dans sa portée spécifique :
construire un assistant adapté aux \emph{supports de cours de première année
du CMI Informatique}, qui présentent des caractéristiques particulières
(informatique et mathématiques de niveau introductif, contenu de projet en
français et en anglais). Les défis que nous anticipons s'alignent avec ceux
rapportés dans la littérature : assurer une couverture documentaire suffisante,
gérer les questions hors périmètre et calibrer la confiance des étudiants
dans les sorties du système.
\newpage

\bibliographystyle{plain}
\bibliography{references}
\end{document}